\subsection{One-Hour-Ahead Forecasting Problem}

Let \(\tau\) denote the forecast origin and
\(\mathcal{F}_{\tau}\) the information available at that time.
The target is next-hour log realized variance,
\(y_{\tau+1}=\log(RV_{\tau+1})\). Each model produces the
conditional forecast

\begin{equation}
\widehat{y}_{\tau+1|\tau}
=
E
\left[
y_{\tau+1}
\mid
\mathcal{F}_{\tau}
\right].
\label{eq:forecast-objective}
\end{equation}

For a forecast issued at 12:00 UTC, the predictors use
information observed by 12:00 UTC and the target is realized
variance from 12:00 to 13:00 UTC.

Lagged variables and rolling averages are constructed on the
complete hourly calendar before model-ready observations are
selected. This preserves the calendar interpretation of each
lag and rolling window when invalid hours occur.

\subsection{Chronological Sample Division}

The sample is divided chronologically into training,
validation, and test periods.

\begin{table}[H]
\centering
\caption{Chronological forecasting design}
\label{tab:forecasting-splits}
\begin{tabular}{lll}
\toprule
Sample & Period & Purpose \\
\midrule
Training
& January 2022--December 2023
& Estimation and preprocessing \\

Validation
& January 2024--December 2024
& Model and hyperparameter selection \\

Test
& January 2025--December 2025
& Final out-of-sample evaluation \\
\bottomrule
\end{tabular}
\end{table}

The training sample determines preprocessing parameters and
initial model estimates. The downside-semivariance scale and
feature-standardization parameters are estimated exclusively
from this period.

The validation sample is used to select the autoregressive
order and evaluate the fixed and time-varying specifications.
After the model structures are fixed, the selected models are
re-estimated using the combined training and validation
sample. The 2025 observations are reserved for the final
evaluation.

\subsection{Common Evaluation Sample}

All models are compared on identical forecast timestamps. An
observation is retained only when the one-hour-ahead target and
all predictors required by HAR-X are available.

The common sample contains \(16{,}673\) training observations,
\(8{,}610\) validation observations, and \(8{,}422\) test
observations.

Persistence and AR(24) could generate forecasts for additional
timestamps. Restricting evaluation to the common sample ensures
that differences in reported losses reflect forecast
performance under the same market conditions.

\subsection{Sequential Forecasting and Information Control}

The fixed-parameter models use predictor values observed at the
forecast origin and retain constant coefficient estimates
throughout each evaluation period.

TVP-HAR-X generates forecasts sequentially. At origin
\(\tau\), the predicted state produces
\(\widehat{y}_{\tau+1|\tau}\). Once \(y_{\tau+1}\) becomes
available, the Kalman filter updates the state distribution for
the following forecast.

Invalid target hours remain on the state-space calendar as
missing observations. The state transition continues through
these hours, while the measurement update is omitted.

\subsection{Model Selection Rules}

The autoregressive order is selected from
\(p\in\{1,2,\ldots,24\}\). Each candidate model is estimated on
the training sample and evaluated on the same validation
timestamps.

Validation QLIKE is the primary selection criterion. Ties are
resolved sequentially using:

\begin{enumerate}[label=(\roman*)]
    \item lower validation RMSE;
    \item lower validation MAE;
    \item lower autoregressive order.
\end{enumerate}

Training-sample AIC and BIC are reported as supplementary
diagnostics. The HAR-X predictor set and TVP-HAR-X state
specification are fixed before the test evaluation.

\subsection{Mean Absolute Error}

Mean absolute error is calculated on the log
realized-variance scale:

\begin{equation}
MAE
=
\frac{1}{N}
\sum_{t=1}^{N}
\left|
y_t-\widehat{y}_t
\right|.
\label{eq:mae}
\end{equation}

MAE reports the average absolute forecast error and gives each
observation a linear contribution to aggregate loss.

\subsection{Root Mean Squared Error}

Root mean squared error is

\begin{equation}
RMSE
=
\sqrt{
\frac{1}{N}
\sum_{t=1}^{N}
\left(
y_t-\widehat{y}_t
\right)^2
}.
\label{eq:rmse}
\end{equation}

The squared-error transformation assigns greater weight to
large forecast errors.

\subsection{QLIKE Loss}

QLIKE is the primary volatility forecast loss. It is suitable
for comparing forecasts when realized variance is an imperfect
proxy for latent conditional variance
\citep{patton2011}.

Let \(RV_t=\exp(y_t)\) and
\(\widehat{RV}_t=\exp(\widehat{y}_t)\). Observation-level
QLIKE is

\begin{equation}
QLIKE_t
=
\frac{RV_t}{\widehat{RV}_t}
-
\log
\left(
\frac{RV_t}{\widehat{RV}_t}
\right)
-1.
\label{eq:qlike-original}
\end{equation}

On the log scale, the same loss can be written as

\begin{equation}
QLIKE_t
=
\exp
\left(
y_t-\widehat{y}_t
\right)
-
\left(
y_t-\widehat{y}_t
\right)
-1.
\label{eq:qlike-log}
\end{equation}

Average QLIKE is

\begin{equation}
QLIKE
=
\frac{1}{N}
\sum_{t=1}^{N}
QLIKE_t.
\label{eq:average-qlike}
\end{equation}

Lower values indicate better performance. QLIKE evaluates
proportional variance errors and can impose a large penalty
when high realized variance is substantially underestimated.

\subsection{Forecast Loss Comparison}

The percentage improvement of model \(m\) over benchmark \(b\)
is

\begin{equation}
Improvement_{m,b}
=
100
\left(
\frac{
L_b-L_m
}{
L_b
}
\right),
\label{eq:percentage-improvement}
\end{equation}

where \(L_b\) and \(L_m\) denote average forecast losses. A
positive value indicates lower loss for model \(m\).

Predictive accuracy is assessed through the loss-differential
framework of \citet{dieboldmariano1995}. For TVP-HAR-X and a
benchmark model \(b\), the loss differential is

\begin{equation}
d_t
=
L_{b,t}
-
L_{\mathrm{TVP},t}.
\label{eq:loss-differential}
\end{equation}

Positive values favor TVP-HAR-X. The hypotheses are

\[
H_0:
E(d_t)\leq 0,
\qquad
H_1:
E(d_t)>0.
\]

Because forecast losses may be heteroskedastic and serially
correlated, the mean loss differential is evaluated using
Newey--West HAC standard errors with 24 lags
\citep{neweywest1987}. The resulting statistic is interpreted
as a Diebold--Mariano-type test of equal predictive accuracy.

The pre-specified primary comparison evaluates TVP-HAR-X
against Fixed HAR-X under QLIKE. Absolute-error and
squared-error comparisons are reported as supplementary
evidence.